\section{Optimal graph couplings for zero transport}
\label{sec:graph-minimax}

The final question of the Article's section \texttt{zer:sec:open} asks which
graph minimizes the largest displacement bound in its electrical zero-transport
theorem, subject to a conductance budget and prescribed input degrees.
We solve this optimization, and its analogue for every Wasserstein distance
$W_p$, $1\le p\le\infty$. The result is unconditional finite-dimensional
mathematics. Dirichlet GRH supplies an application to finite packets of zero
ordinates.

\subsection{The optimization problem and its solution}

Fix $n\ge2$, integers $d_1,\ldots,d_n\ge1$, and a budget $W>0$. Put
\begin{equation}\label{gm:eq:C}
 N=\sum_{i=1}^n d_i,\qquad d=(d_i)_{i=1}^n,\qquad
 C=\operatorname{diag}(d_1,\ldots,d_n)-\frac{dd^{\mathsf T}}N.
\end{equation}
An admissible graph has these $n$ labelled vertices, conductances
$w_{ij}=w_{ji}\ge0$, connected positive support, and
$\sum_{i<j}w_{ij}=W$. Its Laplacian $L$ is positive definite on
$H=\mathbf1_n^\perp$ and has trace $2W$. The topology is free.

At vertex $i$ place a monic real-rooted polynomial $P_i$ of degree $d_i$.
Define
\begin{equation}\label{gm:eq:forest}
 \mathcal F_G(z)=\sum_{F}
 \left(\prod_{e\in F}w_e\right)
 \left(\prod_{i\in\mathcal R(F)}P_i'(z)\right)
 \left(\prod_{i\notin\mathcal R(F)}P_i(z)\right),
\end{equation}
where the sum is over spanning forests with one selected root in each
component and $\mathcal R(F)$ is the set of selected roots. There are $N-1$
real zeros $\lambda_1\le\cdots\le\lambda_{N-1}$ of this polynomial, counted
with multiplicity; a proof is given below. Let
$\beta_1\le\cdots\le\beta_{N-1}$ be the zeros of
$(\prod_iP_i)'$. Form the two empirical probability measures with mass
$1/(N-1)$ at each listed zero. Their one-dimensional Wasserstein distances are
\begin{align}
 W_p(\lambda,\beta)&=
 \left(\frac1{N-1}\sum_{k=1}^{N-1}|\beta_k-\lambda_k|^p\right)^{1/p}
 &&(1\le p<\infty),\label{gm:eq:Wp}\\
 W_\infty(\lambda,\beta)&=\max_k|\beta_k-\lambda_k|.\nonumber
\end{align}
Write $\mathcal E_p(G)$ for the supremum of this distance over all permitted
inputs $P_i$. This is a guarantee based only on their degrees, before their
roots are specified.

\begin{theorem}[Exact minimax graph]\label{gm:thm:main}
For $1\le p<\infty$, set $r=p/(p+1)$ and $S_r=\operatorname{tr}(C^r)$,
with the zero eigenvalue contributing zero. The unique optimal graph has
Laplacian
\begin{equation}\label{gm:eq:optimalL}
 L_p=\frac{2W}{S_r}C^r,
 \qquad
 \min_G\mathcal E_p(G)=
 \frac{S_r^{\,1+1/p}}{2W(N-1)^{1/p}}.
\end{equation}
For the largest displacement, the unique optimum is
\begin{align}
 L_\infty&=\frac{2W}{\operatorname{tr}C}C,
 &w_{ij}^{(\infty)}&=
 \frac{W d_i d_j}{\sum_{a<b}d_a d_b},\label{gm:eq:infinity}\\
 \min_G\mathcal E_\infty(G)&=
 \frac{\operatorname{tr}C}{2W}
 =\frac{N^2-\sum_i d_i^2}{2WN}.\nonumber
\end{align}
Every optimizer is a complete graph with strictly positive conductances.
For each fixed admissible graph, placing every input root at one common real
number attains its supremum. If all input roots are instead required to be
distinct, the supremum and the unique optimal graph are unchanged.
\end{theorem}

At $p=1$ the displacement is independent of the input roots, so no supremum
is needed. For $p>1$ we give an attaining configuration, without claiming that
it is the only attaining configuration. The normalization by $N-1$ in
\eqref{gm:eq:Wp} is essential: the perturbation has rank $n-1$, but there are
$N-1$ output zeros.

\subsection{The compression identity}

We first verify the finite identity behind the Article's electrical theorem.
List all input roots, repeated with their multiplicities and grouped by
vertex, as the diagonal entries of a real matrix $\Lambda$ of size $N$.
Let $V$ be the $N$ by $n$ matrix whose row for a root at vertex $i$ is $e_i^{\mathsf T}$.
Set
\begin{equation}\label{gm:eq:compression-data}
 Q=I_N-\frac{\mathbf1_N\mathbf1_N^{\mathsf T}}N,\quad
 \mathcal V=\mathbf1_N^\perp,\quad
 A=(Q\Lambda Q)|_{\mathcal V},\quad
 B=(QVL^+V^{\mathsf T}Q)|_{\mathcal V},
\end{equation}
where $L^+$ is the Moore--Penrose inverse.

\begin{lemma}[Forest compression]\label{gm:lem:compression}
If $\tau(G)$ is the sum of the conductance products of spanning trees, then
\[
 \mathcal F_G(z)=N\tau(G)\det_{\mathcal V}(zI-A+B),\qquad
 \left(\prod_iP_i(z)\right)'=N\det_{\mathcal V}(zI-A).
\]
Moreover, $B$ is positive semidefinite of rank $n-1$, and its nonzero
eigenvalues are the eigenvalues on $H$ of
\begin{equation}\label{gm:eq:reducedB}
 T_L=C^{1/2}(L|_H)^{-1}C^{1/2}.
\end{equation}
\end{lemma}

\begin{proof}
The matrix-tree theorem, applied after adding one extra vertex with edge
weights $x_i$ to the original vertices, gives the rooted-forest identity
\[
 \det(L+\operatorname{diag}(x_i))
 =\sum_F\left(\prod_{e\in F}w_e\right)
              \prod_{i\in\mathcal R(F)}x_i.
\]
Off the input roots, put $x_i=P_i'/P_i$ and multiply by
$P=\prod_iP_i$. Since
$V^{\mathsf T}(zI-\Lambda)^{-1}V=\operatorname{diag}(P_i'/P_i)$,
the Schur determinant formula gives
\[
 \mathcal F_G(z)=
 \det\begin{pmatrix}zI_N-\Lambda&V\\-V^{\mathsf T}&L\end{pmatrix}.
\]
Both sides are polynomials, so this identity holds everywhere.

Diagonalize $L$ in a basis beginning with
$u=\mathbf1_n/\sqrt n$ and eliminate its positive eigenvalue block.
The product of those eigenvalues is $n\tau(G)$, leaving the bordered
determinant with upper block
$zI_N-\Lambda+VL^+V^{\mathsf T}$ and border
$Vu=\mathbf1_N/\sqrt n$. For any nonzero real vector $b$ and any square
matrix $M$,
\[
 \det\begin{pmatrix}M&b\\-b^{\mathsf T}&0\end{pmatrix}
 =\|b\|^2\det_{b^\perp}(P_{b^\perp}M|_{b^\perp});
\]
expansion in an orthonormal basis ending with $b/\|b\|$ proves this.
Here $\|Vu\|^2=N/n$, proving the first identity.
The second follows from the cofactor identity
\[
 \det_{\mathcal V}(zI-A)
 =P(z)\frac1N\mathbf1_N^{\mathsf T}(zI-\Lambda)^{-1}\mathbf1_N
 =\frac{P'(z)}N.
\]

Finally, $V^{\mathsf T}QV=C$. The equal nonzero spectra of
$XX^{\mathsf T}$ and $X^{\mathsf T}X$, with
$X=QV(L^+)^{1/2}$, identify the positive spectrum of $B$ with that of
$T_L$. The quadratic identity
\[
 y^{\mathsf T}Cy=\frac1N\sum_{i<j}d_id_j(y_i-y_j)^2
\]
shows that $C$ is positive definite on $H$, proving the rank claim.
\end{proof}

In particular the output roots are real. Min--max gives
$\lambda_k\le\beta_k$ for all $k$.

\subsection{An exact worst-case envelope}

The eigenvalue perturbation inequality used next is classical Lidskii
majorization; see, for example, \cite[Theorem~2.1]{argentati2008}.
We include a proof of the specific convex-cost form required here.

\begin{lemma}[Convex displacement bound]\label{gm:lem:convex}
Let $A$ be real symmetric and $B\succeq0$ of the same size $m$.
Write ordered eigenvalues increasingly and let
$\Delta_k=\lambda_k(A)-\lambda_k(A-B)\ge0$.
For every continuous convex $\phi:[0,\infty)\to\mathbb R$,
\[
 \sum_{k=1}^m\phi(\Delta_k)\le\operatorname{tr}\phi(B).
\]
Equality holds for $A=cI$; for $\phi(t)=t$ it holds for every $A$.
\end{lemma}

\begin{proof}
First suppose $A$ has simple spectrum. The discriminant of the characteristic
polynomial of $A-tB$ is a polynomial in $t$ that is nonzero at zero, so only
finitely many $t\in[0,1]$ have repeated eigenvalues. At the other points,
\[
 -\lambda_k'(t)=u_k(t)^{\mathsf T}Bu_k(t)=b_k(t)\ge0
\]
for an orthonormal eigenbasis $u_k(t)$. The ordered eigenvalues are Lipschitz,
so $\Delta_k=\int_0^1b_k(t)\,dt$. Jensen's inequality gives
$\phi(\Delta_k)\le\int_0^1\phi(b_k(t))\,dt$.
In a $B$-eigenbasis, each $b_k(t)$ is a convex combination of its eigenvalues,
with weights equal to squared entries of an orthogonal matrix. A second
application of Jensen, summed over $k$, gives
$\sum_k\phi(b_k(t))\le\operatorname{tr}\phi(B)$.
Integrate. Approximation by matrices with simple spectrum proves the general
case. The two equality assertions follow directly from the spectra or the
trace identity.
\end{proof}

Apply this lemma with $\phi(t)=t^p$, and apply min--max for $p=\infty$.
The choice $P_i(z)=(z-c)^{d_i}$ makes $A=cI$ and attains the bound.
Together with Lemma~\ref{gm:lem:compression}, this proves the exact identities
\begin{equation}\label{gm:eq:exact-objective}
 \mathcal E_p(G)^p=\frac{\operatorname{tr}_H(T_L^p)}{N-1}
 \quad(1\le p<\infty),\qquad
 \mathcal E_\infty(G)=\|T_L\|.
\end{equation}
Arbitrarily small distinct perturbations of a common root give the same
supremum under the distinct-root restriction, by continuity of eigenvalues.

\subsection{Optimization, uniqueness, and graph feasibility}

\begin{proof}[Proof of Theorem~\ref{gm:thm:main}]
Work on $H$ in an orthonormal $C$-eigenbasis, with eigenvalues
$c_1,\ldots,c_{n-1}>0$. Put $\ell_i=L_{ii}>0$, so
$\sum_i\ell_i=2W$. For finite $p\ge1$, spectral Jensen and Cauchy--Schwarz give
\begin{align*}
 \operatorname{tr}(T_L^p)
 &\ge\sum_i ((T_L)_{ii})^p
 =\sum_i\bigl(c_i(L^{-1})_{ii}\bigr)^p
 \ge\sum_i(c_i/\ell_i)^p\\
 &\ge\frac{(\sum_i c_i^{p/(p+1)})^{p+1}}{(2W)^p}.
\end{align*}
Here $(L^{-1})_{ii}\ell_i\ge1$ is Cauchy--Schwarz for
$L^{1/2}e_i$ and $L^{-1/2}e_i$; equality forces $e_i$ to be an
$L$-eigenvector. The last inequality is the scalar constrained minimum of
$\sum_i c_i^p\ell_i^{-p}$, a strictly convex function: its unique minimizer is
$\ell_i=2Wc_i^{p/(p+1)}/S_r$.
Equality throughout therefore forces $L=(2W/S_r)C^r$.
This argument includes $p=1$: its first inequality is equality, but the inverse
diagonal inequality still forces diagonality and the scalar minimum is still
unique.

For $p=\infty$, let $t=\|T_L\|$. The inequality $T_L\preceq tI$
is equivalent to $C\preceq tL$. Thus
$t\ge\operatorname{tr}C/(2W)$. At equality $tL-C$ is positive semidefinite
with zero trace and is zero, giving the unique candidate in
\eqref{gm:eq:infinity}.

These lower bounds hold for all positive definite $L|_H$ of the prescribed
trace, even without the sign restrictions of a graph Laplacian. It remains
to verify that the candidates are graphs. For $0<r<1$, the scalar beta
integral and spectral calculus give
\[
 C^r=\frac{\sin(\pi r)}\pi\int_0^\infty
 t^{r-1}C(C+tI)^{-1}\,dt.
\]
This matrix annihilates $\mathbf1_n$. For $i\ne j$, the rank-one inverse
formula applied to $C+tI=\operatorname{diag}(d_i+t)-dd^{\mathsf T}/N$ gives
\[
 ((C+tI)^{-1})_{ij}=
 \frac{[d_i/(t+d_i)][d_j/(t+d_j)]}
 {t\sum_\ell d_\ell/(t+d_\ell)}>0.
\]
Consequently
\begin{equation}\label{gm:eq:edge-integral}
 -(C^r)_{ij}=\frac{\sin(\pi r)}\pi
 \int_0^\infty
 \frac{t^{r-1}d_id_j}
 {(t+d_i)(t+d_j)\sum_\ell d_\ell/(t+d_\ell)}\,dt>0.
\end{equation}
The integrand is $O(t^{r-1})$ at zero and $O(t^{r-2})$ at infinity.
Thus $C^r$ has strictly negative off-diagonal entries and is a connected
complete-graph Laplacian. At $r=1$ this follows from
$C_{ij}=-d_id_j/N$. Scaling gives trace $2W$ and the required budget.
Equation~\eqref{gm:eq:exact-objective} now gives the stated minima.
\end{proof}

Fractional powers of graph Laplacians are established objects;
\cite[Theorem~3.1]{benzi2022} gives a general preservation theorem.
Formula~\eqref{gm:eq:edge-integral} proves feasibility here directly and
provides an explicit positive integral for every optimal edge weight.

\begin{corollary}[Dependence on the displacement cost]\label{gm:cor:different}
If $n\ge3$, two distinct exponents $p\in[1,\infty]$ yield the same optimal
graph exactly when $d_1=\cdots=d_n$. In that case every optimal edge has
conductance $2W/[n(n-1)]$. For $n=2$ all exponents give the sole available
edge, of conductance $W$.
\end{corollary}

\begin{proof}
Equality of two normalized powers of $C$ with different exponents forces all
its positive eigenvalues to coincide. Thus $C=c(I_n-\mathbf1_n\mathbf1_n^{\mathsf T}/n)$.
Its off-diagonal entries imply that all products $d_id_j$ for $i\ne j$ are
equal; for $n\ge3$ this is equivalent to equal degrees. The converse and the
two-vertex case follow immediately.
\end{proof}

\subsection{An exact certificate for near-optimal mean transport}

\begin{proposition}[Mean-error gap and stability]\label{gm:prop:stability}
Let $L_*=L_1$, let $L$ be any admissible Laplacian, and set
$\Delta=L-L_*$ and $\delta=\mathcal E_1(L)/\mathcal E_1(L_*)-1$.
Then
\begin{equation}\label{gm:eq:gap}
 \delta=\frac1{2W}\operatorname{tr}(\Delta L^+\Delta),\qquad
 \|L-L_*\|_{\mathrm F}\le2W\sqrt\delta.
\end{equation}
For $n\ge3$, the exponent $1/2$ cannot be improved in a local bound uniform
over admissible Laplacians.
\end{proposition}

\begin{proof}
All inverses in this proof act on $H$. Write
$\alpha=2W/\operatorname{tr}\sqrt C$, so $C=L_*^2/\alpha^2$.
The exact resolvent identity is
\[
 L^{-1}=L_*^{-1}-L_*^{-1}\Delta L_*^{-1}
       +L_*^{-1}\Delta L^{-1}\Delta L_*^{-1}.
\]
Multiply by $C$ and take traces. Since $\operatorname{tr}\Delta=0$,
\[
 \operatorname{tr}(CL^{-1})-\operatorname{tr}(CL_*^{-1})
 =\alpha^{-2}\operatorname{tr}(\Delta L^{-1}\Delta),\qquad
 \operatorname{tr}(CL_*^{-1})=2W/\alpha^2.
\]
Equation~\eqref{gm:eq:exact-objective} proves the first claim.
Since $\lambda_{\max}(L)\le\operatorname{tr}L=2W$, the trace on the
right is at least $\|\Delta\|_{\mathrm F}^2/(2W)$.
For sharpness of the exponent, exchange a small amount of conductance
between two edges of the strictly positive optimal complete graph. This is
$L(t)=L_*+tK$ for nonzero symmetric $K$ with $K\mathbf1_n=0$
and $\operatorname{tr}K=0$. The exact formula gives
$\delta(t)=t^2\operatorname{tr}(K L_*^+K)/(2W)+O(t^3)$ with positive
quadratic coefficient, whereas $\|L(t)-L_*\|_{\mathrm F}=|t|\|K\|_{\mathrm F}$.
\end{proof}

\subsection{Unequal degrees: average and worst displacement disagree}

Take $(d_1,d_2,d_3)=(1,1,2)$ and $W=1$. The positive eigenvalues of $C$
are $1$ and $3/2$, with eigenvectors $(1,-1,0)$ and $(1,1,-2)$.
For $r=p/(p+1)$ and $s=(3/2)^r$, formula~\eqref{gm:eq:optimalL} gives
\[
 w_{12}=\frac{3-s}{3(1+s)},\qquad
 w_{13}=w_{23}=\frac{2s}{3(1+s)}.
\]
At $p=\infty$ these are $(1/5,2/5,2/5)$ and the optimal maximum
displacement is $5/4$. At $p=1$, $s=\sqrt6/2$, the edges are
\[
 \left(-3+\frac{4\sqrt6}3,\quad
 2-\frac{2\sqrt6}3,\quad 2-\frac{2\sqrt6}3\right),
\]
and the optimal mean displacement is $(5+2\sqrt6)/12$.
The graph optimal for the maximum has mean $5/6$, strictly larger than
$(5+2\sqrt6)/12$; the graph optimal for the mean has maximum
$(3+\sqrt6)/4$, strictly larger than $5/4$.

For the attaining inputs $(P_1,P_2,P_3)=(z,z,z^2)$ and conductances
$(w_{12},w_{13},w_{23})=(x,y,y)$, direct expansion gives
\[
 \mathcal F_G(z)=2z+(4x+6y)z^2+4(2xy+y^2)z^3.
\]
At the maximum-optimal graph this factors as
$\frac{32}{25}z(z+\frac54)^2$, whereas $(P_1P_2P_3)'=4z^3$.
At the mean-optimal graph the displacement multiset is
$\{0,(1+s)/2,s(1+s)/2\}$.

The supplied verifier checks these identities symbolically, an exact
noncollapsed forest/compression example, and the stability certificate.
Its numerical diagnostics also check $75$ optimal configurations,
$750$ competing graphs and $150$ checks comparing noncollapsed inputs
with the collapsed extremizer. Those diagnostics check implementation and constants;
the theorem rests on the proof above.

\subsection{Application to finite Dirichlet-zero packets}

Under the Dirichlet GRH theorem in the vendored development, the phase-normalized
completed primitive Dirichlet $L$-functions have real ordinate zeros. Select a
finite packet of $d_i$ such zeros at vertex $i$, preserving the selected
multiplicities, and make $P_i$ their monic root polynomial. Theorem~\ref{gm:thm:main}
gives the best uniform guarantee based only on the packet sizes and budget.
It allows shared zeros, multiple zeros and arbitrary spacings.
A particular arithmetic packet may admit a smaller error and a different
optimal graph; the common-root extremizer need not occur for Dirichlet
functions. No optimum for an entire infinite zero set is inferred.
